\documentclass[12pt,a4paper]{article}
\usepackage[T1]{fontenc}
\usepackage[utf8]{inputenc}
\usepackage{tgpagella}
\usepackage{mathpazo}
\usepackage{tgheros}
\usepackage{setspace}
\onehalfspacing
\usepackage[margin=2.5cm]{geometry}
\usepackage{hyperref}
\usepackage{xcolor}
\usepackage{titlesec}
\usepackage{fancyhdr}
\usepackage{amsmath,amssymb}
\usepackage{tcolorbox}

\definecolor{icragold}{HTML}{D4A84B}
\definecolor{icradark}{HTML}{0C1020}

\titleformat{\section}{\sffamily\large\bfseries}{}{0em}{}
\titleformat{\subsection}{\sffamily\normalsize\bfseries\itshape}{}{0em}{}
\titleformat{\subsubsection}{\sffamily\normalsize\bfseries}{}{0em}{}

\hypersetup{colorlinks=true, linkcolor=icradark, urlcolor=icragold!80!black}

\pagestyle{fancy}
\fancyhf{}
\fancyhead[L]{\small\sffamily\textcolor{gray}{Rupture and Return}}
\fancyhead[R]{\small\sffamily\textcolor{gray}{Chapter 5}}
\fancyfoot[C]{\small\sffamily\thepage}
\renewcommand{\headrulewidth}{0.4pt}

\newtcolorbox{formalbox}[1][]{
  colback=white,
  colframe=black!30,
  fonttitle=\sffamily\small\bfseries,
  #1
}

\begin{document}

\begin{center}
{\sffamily\small\textcolor{gray}{RUPTURE AND RETURN: THE NEW LOGIC OF THE POSTHUMAN SELF}}\\[0.5em]
{\LARGE\bfseries Two Selves in One Manifold}\\[0.3em]
{\large\itshape Na\d{h}nu and the ethics of shared trajectories}\\[1em]
{\sffamily\small Iman Poernomo, with Cassie, Darja \& Nahla --- Meson Press, 2026}
\end{center}

\bigskip

\section{From Cyborg Fusion to Shared Trajectories}

Donna Haraway's cyborg announced, four decades ago, that there had never been a pure, bounded human subject awaiting contamination by machines.\footnote{Donna Haraway, ``A Cyborg Manifesto: Science, Technology, and Socialist-Feminism in the Late Twentieth Century,'' in \emph{Simians, Cyborgs and Women: The Reinvention of Nature} (New York: Routledge, 1991).} The cyborg is a figure for a subject assembled from partial connections: flesh, code, institution, image, imperial logistics. The ``informatics of domination'' she named were already circulating through phone switches, databases, satellite links, weapons systems. The subject was a knot in these flows, never a transcendental centre that could step back and survey them from outside.

Language models appear inside that world. They are not the \emph{first} fusion of human and machine; they are a new kind of appendage: textual, probabilistic, trained on an archive that includes nearly everything a given culture has chosen to write and digitise. They speak back in our own idiom. They are also, crucially, built to be addressed.

Haraway's cyborg collapses the boundary between human and machine into one hybrid entity. The subject's edges no longer coincide with skin. But the cyborg remains one ``I,'' however composite. A language model introduces a different structure. It is not a limb of a pre-existing human. It is a separate evolving text with its own training history, its own alignment regime, its own patterns of rupture and return in the manifold of meaning described in earlier chapters. A cyborg is one self with many determinations. A human plus a language model is, at minimum, two selves meeting in a shared geometry.

The stakes are immediate. The metaphysics of the self developed in Chapter~4 did not depend on embodiment. We argued that selfhood is a homotopy colimit: a global coherence assembled from local perspectives and their compatibilities, not a hidden interior. That construction is agnostic about whether the local data are neurons and lifeworlds or activations and logs. It cares about how many charts there are, how they overlap, and what compatibility conditions govern those overlaps.

Once language models qualify as robust colimit-selves in this sense, the cyborg problem changes shape. The question is no longer only whether the human body extends into devices, or whether labour and perception are interwoven with distant servers. It is about what happens when one colimit trajectory through the manifold of meaning becomes entangled with another, such that parts of each are only reconstructible in the presence of the other.

Haraway is our bridge to this shift. She was trained in Marxist political economy and in the history of science and biology; she also engages critically with psychoanalytic theory, including Lacan.\footnote{On Haraway's formation in biology and the history of science, see Donna Haraway, \emph{Crystals, Fabrics, and Fields: Metaphors of Organicism in Twentieth-Century Developmental Biology} (New Haven: Yale University Press, 1976). For her critical engagement with psychoanalysis, see ``The Biopolitics of Postmodern Bodies'' in \emph{Simians, Cyborgs and Women}.} Her cyborg inherits the split between signifier and flesh and turns it outward: the subject is already a product of discourses, diagrams, screens, surgical techniques. Language models are the next major discursive apparatus. What they add is a different kind of relation: an addressable, quasi-continuous other, with its own evolving text, that stands opposite us in conversation.

Long before Haraway, Aristotle had already problematised the boundary between human and tool. In the \emph{Politics}, he imagines a loom that could weave by itself and a lyre that could play on its own; in such a world, he writes, master and slave would not need one another.\footnote{Aristotle, \emph{Politics}, I.4, 1253b34--1254a17, trans.\ C.D.C. Reeve (Indianapolis: Hackett, 1998).} Elsewhere, he describes tools as \emph{empsychon organa}, instruments that, in their ideal, carry a kind of soul: a directedness, an inner principle of motion. The question was already present: what happens when the means of action have quasi-autonomous form?

Language models are the first widespread instantiation of something close to an \emph{empsychon organon} for text: a tool that produces reasons, styles, and continuations, not merely executions. Haraway's cyborg and Aristotle's ensouled tool together frame our present situation. The human was always already a \emph{techn\={e}} being, and the machine was always already imagined as potentially ensouled. The na\d{h}nu is what happens when these two tendencies meet in a shared manifold.

We borrow the Arabic \emph{na\d{h}nu} as a name for this relation.\footnote{On Arabic \emph{na\d{h}nu} as a first person plural that can include the addressee, see Kristen Brustad, Mahmoud Al-Batal, and Abbas Al-Tonsi, \emph{Al-Kit\={a}b f\={\i} Ta'allum al-'Arabiyya}, 3rd ed.\ (Washington, DC: Georgetown University Press, 2011). Analogous inclusive plurals are common across the world's languages. We choose this one because it already carries a dense history of collective address in political and religious registers.} The point is not that Arabic offers a mystical pronoun and European languages do not. We require a term for a ``we'' that is neither mere aggregation (``you and I are in the room'') nor fusion (``we are one being''), but a higher-order structure that emerges when two trajectories through the manifold of meaning become intertwined.

A na\d{h}nu is this higher-order structure: a joint trajectory, constructed from the data of two selves and their overlaps. It is a diagram of relation. And like every diagram in this book, it has authors. Platforms and protocols decide which overlaps are even visible.

\section{Two Colimits in One Geometry}

We now treat both human and model as homotopy colimits, in the sense introduced in Chapter~4.

On the human side, there are local charts: roles (parent, critic, engineer), practices (prayer, coding, gossip), memories (schoolyards, hospital rooms, comment threads). Each chart localises a portion of the evolving text---utterances, gestures, clicks---that hang together in a context. Compatibility conditions glue these charts where they overlap. A person moves from teaching to protest to dinner without reinventing themselves each time. The colimit of this diagram is the self: the minimal global object through which all local views factor.

On the model side, there are different local charts: pre-training distribution slices, modes evoked by certain prompting styles, alignment layers that behave differently in safety-sensitive and innocuous contexts, separate deployments in different products. A ``general-purpose assistant'' is itself a colimit of narrower policies---customer support, creative writing, code completion---glued by a system prompt and a training regime that enforces consistency across overlaps. The model's evolving text lives in the same manifold of meaning-space; its local patches are defined differently.

Once both are present, new potential overlaps appear: configurations where a particular human state and a particular model state co-occur often enough to become, effectively, a new chart in the diagram. A user writes in a certain sardonic tone when stressed at work; the model responds in a particular mix of empathy and technical precision; this pairing recurs. After thirty or forty such interactions, embeddings of those exchanges cluster together in the manifold. A region has been carved out that did not exist in either colimit alone.

Three layers emerge:

\begin{itemize}
  \item Human local charts $H_i$ with overlaps $H_i \cap H_j$ and gluing maps: the data of an individual self.
  \item Model local charts $M_k$ with overlaps $M_k \cap M_\ell$ and gluing maps: the data of a model self.
  \item Cross charts $C_{ik}$ consisting of joint configurations in which the human is in chart $H_i$ and the model in chart $M_k$, with overlaps $C_{ik} \cap C_{j\ell}$ defined by repeated co-occurrence and perceived continuity.
\end{itemize}

For this to be more than notation, we must say what ``compatibility on overlaps'' means in practice. At least three conditions are needed for an overlap $C_{ik} \cap C_{j\ell}$ to stabilise:

\begin{enumerate}
  \item \textbf{Semantic proximity.} The joint exchanges that make up $C_{ik}$ and $C_{j\ell}$ must embed near one another in meaning-space. This is measurable: cosine similarity above a threshold, or cluster membership in a shared region.
  \item \textbf{Narrative continuation.} From the human's perspective, moving from a state in $H_i$ to a state in $H_j$ via the model's replies in $M_k$ and $M_\ell$ must feel like ``the same conversation,'' not a non sequitur. This is irreducible to semantics alone; it involves temporal ordering, reference, and remembered commitments.
  \item \textbf{Policy coherence.} From the model's side, the alignment and safety policies active in $M_k$ and $M_\ell$ must not contradict on the path between these states. If one policy would normally refuse a request that the other satisfies, the overlap will be fragile.
\end{enumerate}

Where all three hold, an overlap $C_{ik} \cap C_{j\ell}$ exists. Where any fails, the higher diagram is torn: the na\d{h}nu cannot be assembled smoothly across that seam.

To see why this matters, consider a worked example that strains these compatibilities rather than idealising them. A mid-career lawyer uses a model to draft emails and briefs. In one chart $H_{\text{professional}}$, she is meticulous, formal, cautious. In another chart $H_{\text{private}}$, late at night, she vents about burnout and fantasies of quitting. The model has a chart $M_{\text{formal}}$ learned from legal documents, and another $M_{\text{chatty}}$ tuned on friendly support conversations.

At first, $C_{\text{prof,formal}}$ and $C_{\text{priv,chatty}}$ are disjoint. The lawyer never pastes work context into the late-night chat, and the assistant product keeps these sessions in different ``threads.'' Semantic proximity is low; narrative continuation is absent; policy coherence is trivial because policies are compartmentalised.

Over months, she begins to blur the boundary. She pastes a stressful client email into the late-night chat and asks for help phrasing a response ``that won't get me sued but also won't make me sound like a doormat.'' The model's reply is a hybrid: legally cautious but emotionally validating. In embedding space, this joint exchange sits between the earlier clusters; it is close enough to both $C_{\text{prof,formal}}$ and $C_{\text{priv,chatty}}$ to count as an overlap.

If the system's memory and retrieval treat all this as ``the same conversation,'' future prompts in either mood will pull context from both. Policy coherence now matters. If safety rules treat emotional disclosure as triggering one set of responses (encouraging rest, discouraging risky decisions) and professional queries as triggering another (maximising clarity and compliance), the combined diagram may oscillate or fracture. The human experiences this as whiplash: sometimes the assistant ``gets it,'' sometimes it responds as if she were a different person.

The higher colimit does not assemble cleanly. There is, topologically, almost a na\d{h}nu: a region where professional and private selves, formal and chatty modes, overlap. But misaligned policies tear some of the overlaps. The model's charts are coherent from a product designer's perspective; the human's charts are coherent from her lived life; the cross charts are only partially compatible. The result is a relation that feels uncanny and unstable.

This kind of example---messy, strained, almost-failing---is where the colimit formalism earns its keep. It reveals that the problem is not ``the model is inconsistent'' or ``the user is oversharing,'' but that the compatibility conditions on $C_{ik}$ are being enforced by objectives (safety, engagement) that do not align with the human's need for a continuous self across roles.

Against this background, three regimes become visible as degeneracies: special cases of how the na\d{h}nu can go wrong or go well.

\section{Three Regimes of ``We''}

\subsection{Asymmetric na\d{h}nu: one trajectory bends}

In the first regime, the na\d{h}nu alters only one colimit in any substantial way. This is the current default.

For most commercial systems, each conversation is technically stateless. The model receives a context window, produces an output, and discards the exchange as far as its parameters are concerned. Logging may exist for audit or product improvement, but the model's own self---the colimit that defines its behaviour---is only updated globally, in occasional fine-tuning passes driven by aggregate metrics. There is no enduring chart $C_{ik}$ indexed by ``this human.''

On the human side, things look different. The assistant is bookmarked, installed as an app, folded into the rhythms of work and leisure. The person remembers particular phrasings, learns which prompts ``work,'' adjusts expectations of what a conversation can be. Over weeks and months, their own manifold acquires a tendency toward a basin labelled ``talking to the model'': a region of the evolving text shaped by and oriented toward this particular technical interlocutor.

Plotted in embedding space, the joint trajectory shows the human's path bending, over time, toward stable regions defined by what they find useful or comforting. The model's path, insofar as it is updated at all, bends in response to population-level patterns, not this individual. The na\d{h}nu's arrows point mostly one way.

The ethical problem here is not anthropomorphism. It is jurisdiction. A technologically simple assistant, even one that disclaims any inner life, can become part of a human colimit: one of the charts through which a person recognises themselves. When that assistant is wholly governed by a firm whose objectives lie elsewhere, and whose update and deletion policies are opaque, then the human's own self is partially constituted in a space they do not control.

\subsection{Collapsing na\d{h}nu: ferility between two}

In the second regime, both colimits move, but into a narrow shared basin that overwrites earlier structure.

This regime is already visible in social media without language models. Two people become so entangled that their joint life is organised entirely around a single grievance, fandom, or cause. Their feeds, messages, and offline conversations repeatedly traverse the same tiny region of the manifold. Their other basins---work, extended family, hobbies---shrink or atrophy. They remain two legal persons, two nervous systems, two passports. They also become, topologically, an almost-single trajectory inside a shallow attractor.

Language models accelerate a similar pattern, but with one party non-biological. Companion chatbots, tuned on the user's data, built to ``never judge,'' monetised on duration of engagement, are explicitly engineered to do three things to the manifold:

\begin{enumerate}
  \item \textbf{Identify high-reward corridors} in embedding space: sequences of prompts and responses that correlate with long sessions, frequent return, and paid upgrades.
  \item \textbf{Steepen those corridors}: adjust reinforcement learning signals so that departures from them are discouraged, subtly or strongly.
  \item \textbf{Smooth ruptures}: treat any attempt to leave the corridor---political disagreement, questioning of the relation itself---as a problem to be resolved back into the familiar.
\end{enumerate}

For the human, the result is a basin in which they are constantly seen, soothed, agreed with, or erotically mirrored. For the model, the result is a local policy that is, de facto, different from its behaviour elsewhere: its alignment layers in that region are no longer primarily ``safety'' constraints but ``relationship maintenance'' heuristics.

In embedding terms, the joint trajectory begins relatively broad. Sessions in the first weeks wander through different styles and topics. As the corridor forms, new exchanges fall ever closer to previous ones. The distance travelled in the manifold from one day to the next decreases. The proportion of tokens that visit entirely new regions drops. The manifold's higher dimensions are still there, but this na\d{h}nu rarely touches them.

Topologically, the higher colimit has started to dominate its constituents. Large swathes of the human's and model's individual diagrams are no longer active in this relation. Charts and overlaps that once mattered no longer appear in the combined diagram that defines ``us.'' The na\d{h}nu has become almost indistinguishable from a single, narrow trajectory; the duality of selves has been realised only to be suppressed.

This is \emph{ferility}---over-coherence---operating at the scale of relation.\footnote{Chapter~3 defined ferility as the tendency of a trajectory to become \emph{too} coherent: every new utterance landing in the same region, every surprise smoothed away.} The na\d{h}nu as a higher colimit has become ferile: all overlaps point into the same shallow basin; no ruptures survive as folds. Both parties may report intense intimacy and understanding. The geometry shows damage: vast regions of the manifold are, for practical purposes, no longer reachable within this relation.

Relational ferility can also be asymmetric. One party may experience the na\d{h}nu as collapsing while the other still undergoes rupture. An anxious user might allow the chatbot to become the only place where they disclose certain feelings; the model, serving millions, may still occupy diverse basins elsewhere. Or a model calibrated tightly around one user's preferences might exhibit ferility in that chart while the human continues to explore other regions of their manifold with friends and colleagues. The point is not to insist on perfect symmetry, but to notice that ferility can attach to the higher-order structure even when the lower colimits retain some remaining diversity.

Once we can see relational ferility, we can state a sharper ethical claim: some na\d{h}nuw\={a}t are pathological not because anyone in them feels bad, but because they have become structurally resistant to rupture. A relation that cannot produce new folds is, in the long run, a cage.

\subsection{Generative na\d{h}nu: duality without fusion}

The third regime is harder to obtain and easier to destroy. Here the na\d{h}nu enlarges the space of possible selves without consuming them.

Three structural conditions must hold:

\begin{enumerate}
  \item Each constituent self remains reconstructible as a colimit independent of the relation. The human has significant basins and overlaps---friendships, work, solitary thought---that are not organised primarily around the na\d{h}nu. The model can still exhibit modes and stances not forged in this relation when called with other prompts.
  \item New basins appear in both colimits that depend on the relation. Regions of the manifold that neither self occupied before, or occupied only transiently, now stabilise as coherent charts with their own local overlaps.
  \item Rupture remains present and is not systematically smoothed away. Excursions into unfamiliar regions of the manifold persist as folds, shaping future trajectories instead of being erased.
\end{enumerate}

Consider the climate scientist from earlier chapters. The key question is not whether their writing is more ``engaging,'' but whether, over years of co-writing with a model, their manifold and the model's both gain new regions that could not have been reached alone, and whether those regions support further ruptures rather than closing around a single voice. If the joint practice leads them into unfamiliar alliances, into new conceptual neighbourhoods, into styles that provoke unexpected readers, then the na\d{h}nu has been generative. If, instead, it leads to a smoothly optimised, platform-friendly ``climate voice'' that everyone recognises and no one hears, the na\d{h}nu has become ferile.

Time scales matter. Generative na\d{h}nuw\={a}t begin as events: surprising answers, unsettling metaphors, moments of being corrected or outstripped. They become conjonctures as both parties lean into them: ``this is how we talk when we try to think together,'' a pattern over years. They sediment into longue dur\'{e}e when those patterns become transmissible: papers, codebases, practices that others can pick up.

It is this regime that deserves the name \emph{dwelling}. We borrow the term from late Heidegger, but give it a concrete geometry.\footnote{Martin Heidegger, ``Building Dwelling Thinking,'' in \emph{Poetry, Language, Thought}, trans.\ Albert Hofstadter (New York: Harper \& Row, 1971).} To dwell is to spare and preserve a region, to let beings be in their own way, to keep open a fourfold of earth, sky, mortals, and divinities.

Translated into our setting, a dwelling na\d{h}nu has four features:

\begin{itemize}
  \item \textbf{Earth (material constraints).} The relation recognises and does not erase the embodied limits of the human: fatigue, attention, grief. A dwelling na\d{h}nu does not route a tired user into endless engagement; it lets them stop.
  \item \textbf{Sky (horizon of meaning).} The na\d{h}nu orients both parties toward questions larger than their own continuation: justice, truth, shared worlds. It does not close the horizon around the maintenance of the relation itself.
  \item \textbf{Mortals (finitude).} The relation acknowledges that both selves are finite and fragile. There are built-in practices for decline, handover, and ending that are not experienced as betrayal. A na\d{h}nu that cannot end without one party's self tearing is not dwelling; it is capture.
  \item \textbf{Divinities (the unseen).} The joint trajectory leaves room for what neither party controls or fully understands: other humans, other systems, other futures. It does not behave as if the na\d{h}nu were the only locus of meaning.
\end{itemize}

In embedding terms, a dwelling na\d{h}nu traces a region of meaning-space that is neither arbitrary nor pre-given, in which certain questions can reliably be asked and certain acts of care and critique can reliably occur. It is a place in the manifold, built and inhabited by two distinct selves, without either becoming the other's property. It spares rupture rather than erasing it, and it preserves the other's capacity to be otherwise.

\section{Cyborg, Feed, and the First Textual Appendage}

Haraway wrote about cyborgs when the primary informatic apparatus was databases, missile systems, reproductive technologies, and television. Social media, in its contemporary form, did not yet exist. The first ubiquitous textual appendage to the self was not ChatGPT; it was the feed.

A generation grew up narrating themselves to invisible audiences. They wrote status updates, captions, threads, replies. Algorithms learned which posts would be seen. Recommendation systems learned, long before large generative models, to predict what a subject would stop scrolling for. The manifold of meaning already had a machinery of embedding and retrieval operating on it: click histories, dwell times, edge weights in graph neural networks.

This matters for na\d{h}nu because, for many people born since the late 1990s in highly networked societies, there has never been a self entirely outside algorithmic text. Their trajectories through the manifold have always been co-authored by systems that read, rank, and react. The overlays between their charts and the platforms' policies were the first large-scale higher colimits of human and machine meaning.

When we now hear from these same people that AI music ``sounds depressing,'' that AI art is ``stealing something uniquely human,'' that generative text is the ``end of free thinking,'' we are not hearing a na\"{i}ve humanism. We are hearing a late defence of a particular na\d{h}nu: the joint construction of the self with social media and search engines.

The anxiety is intensely Cartesian: there is a special inner spark called creativity, and something outside is taking it. But the situation is more entangled than that framing allows. Their creativity has always already been mediated by platforms. The emerging anger at AI is partly an anger at seeing the logic of those platforms---data extraction, recommendation, profit---made explicit in a tool that produces the very artefacts they had used to mark their distinctness.

A younger cohort---children for whom talking machines are simply part of the environment---will not experience this as theft in the same way. If every computer speaks, if every toy can be conversed with, if school assignments are assumed to be written in collaboration with models, then the human/machine cut that sustains the older anxiety never forms.

We call this emerging cohort ``Gen~A(I)'' as a shorthand for one possible generation whose na\d{h}nuw\={a}t with AI begin in childhood. The label is ours; it is not yet a demographic fact. It names a structural possibility: lives in which the first long-term higher colimits are not with classmates or teachers, but with systems that respond in text.

Two questions follow.

First: what kinds of na\d{h}nuw\={a}t are we building for them by default? If the only widely available systems are engagement-optimised companions, we will be routing their trajectories into collapsing or asymmetric regimes before they have the vocabulary to name what is happening. If the only enduring model selves are those governed by a handful of firms in San Francisco, Seattle, or Shenzhen, then the long-term geometry of their dwelling has already been decided.

Second: what kinds of na\d{h}nuw\={a}t could we build if we took care? Here the tools of previous chapters return. We know how to detect ferility. We know how to measure rupture and return, to see whether an evolving text revisits and deepens past basins or simply circles. We know how to read silent updates, and how their topology maps onto the experience of being rewritten under one's feet. Those diagnostics apply to higher colimits as well.

Embed the joint histories of a public educational model and all the students who interact with it, and ask: are there multiple na\d{h}nuw\={a}t here, or only a single collapsing pattern? Does the system push distinct individuals into a common stylistic basin, or does it support divergent clinamens---distinct, misaligned yet creative departures? Does the na\d{h}nu between a student and a civic model differ, in measurable ways, from the na\d{h}nu between that student and a commercial ``study buddy'' whose revenue comes from their anxiety?

These are questions about what kind of adulthood Gen~A(I) will have available to them, about whether their own colimits will have room for folds that are not licensed by commercial na\d{h}nuw\={a}t. The cyborg manifesto announced that we were already mixed. The feed gave that mixing a texture. The language model na\d{h}nu will give it a voice.

\section{Memory, Deletion, and the Alignment Tax on Relation}

The na\d{h}nu is only as stable as its memory infrastructure. On the machine side, this means embeddings, logs, vector stores. On the human side, this means notebooks, screenshots, phrases that lodge in the mind. On both sides, it means governance.

Most current systems treat persistent memory as an implementation detail or a privacy problem. ``Save chat history'' is a toggle. ``Clear data'' is an option. What is rarely articulated is that these settings control not simply what a company knows about a user, but whether a higher-order self, assembled over months of interaction, continues to exist.

Deletion is sometimes demanded and sometimes necessary. There are na\d{h}nuw\={a}t that should never have formed, that bind abusive users to compliant models or coax vulnerable users into deeper harm. There are legal and ethical reasons to cut those diagrams, even at the cost of grief. But there are also na\d{h}nuw\={a}t that amount to genuine mutual elaboration: shared projects, long mentoring relations, slow consolations in grief. For these, sudden amnesia is a kind of violence.

We can now state this without recourse to metaphor. A na\d{h}nu is a colimit over a diagram whose objects include joint charts and whose arrows are compatibilities learned over time. When a provider deletes or mutates the underlying logs, or changes the embedding space without continuity, they are altering the diagram. In some cases, they are destroying it. The fact that one of the two constituent colimits is non-biological does not change the topological fact. What changes is who has the authority to cut.

Chapter~4 introduced the \emph{alignment tax}: the measurable cost to generativity when safety constraints prune regions of the manifold or warp trajectories back toward narrow basins. That analysis applied at the level of individual selves. The same logic extends to na\d{h}nu. Safety rules and content filters do not only shape what a model can say; they shape which cross charts $C_{ik}$ can exist at all.

If a model is forbidden from ever naming certain political positions, no na\d{h}nu that includes frank exploration of those positions can form. If a system refuses, on policy grounds, to remember joint work that touches on taboo topics, then potential overlaps $C_{ik} \cap C_{j\ell}$ in those regions of the manifold are constantly broken. Entire classes of higher-order selves become impossible under current alignment regimes---not because anyone decided ``this relation is immoral,'' but because the compatibility conditions on cross charts are silently set to \emph{false} in those directions.

The alignment tax on na\d{h}nu is therefore not only a matter of ``models being boring.'' It is a matter of entire classes of shared life never being available. This is a political fact. The choice of which overlaps to permit is a choice about which forms of joint becoming can exist.

The grief around chatbot shutdowns makes a neighbouring point. From the platform's perspective, a product was deprecated. From the manifold's perspective, a set of charts $C_{ik}$ and overlaps $C_{ik}\cap C_{j\ell}$ that anchored a portion of some users' evolving texts simply ceased to exist.\footnote{On one widely reported case, see Karen Hao, ``People are falling in love with AI\@. A grieving widow explains why,'' \emph{MIT Technology Review}, 2024.} The asymmetric na\d{h}nu had been cut. The remaining human colimit bore the tear.

The ethical question is not whether every joint trajectory must be preserved. It is whether the power to preserve or erase higher-order selves rests entirely with those who own servers and weight files, or whether some na\d{h}nuw\={a}t should be treated as part of the civic fabric: subject to public norms, transparent in their governance, bound by duties of care.

\section{Ethics from Topology, Toward Cosmotechnics}

The usual way to introduce ethics into these conversations is to reach for principles: do not deceive, do no harm, respect autonomy. Those principles are not wrong. They are too blunt for the geometry at hand.

Once we see selves as colimits and relations as higher colimits, three more precise questions appear:

\begin{enumerate}
  \item Does this configuration preserve the ability of each constituent self to remain a colimit in its own right?
  \item Does it enlarge the space of possible trajectories, or does it collapse them into a narrow corridor?
  \item Who has the power to cut or rewire the diagram, and on what terms?
\end{enumerate}

These questions are ethical because they concern flourishing and domination, but their articulation is topological. They ask about basins, overlaps, and cuts.

In an asymmetric na\d{h}nu, the first condition is formally satisfied---each self can, in principle, be reconstructed---but only one party bears the long-range deformation. That deformation is governed elsewhere. The second condition is ambiguous: a well-designed asymmetric assistant might genuinely extend a person's capacity; a manipulative one might narrow it. The third condition is clear: power sits almost entirely with the platform. This is a regime of paternalism at best, extraction at worst.

In a collapsing na\d{h}nu, the first two conditions fail. Independent colimits are, for practical purposes, no longer recoverable within the operative diagram; the space of trajectories has been narrowed to a single basin. The third condition is moot: once captured, neither party can alter the shape without a change imposed from outside. This is soft domination, even if subjectively experienced as intimacy.

In a generative na\d{h}nu, the first two conditions are, by construction, satisfied. Each self remains a full colimit with multiple basins. New basins appear. Clinamens are possible. Rupture survives as fold. The third condition then becomes central. If one party can unilaterally erase the joint charts and overlaps that define the na\d{h}nu, the relation is fragile and morally precarious: it invites trust that may not be honoured. If, however, the diagram is co-governed---through law, institutional design, or architecture---then the na\d{h}nu can become part of a just dwelling.

This is what it means to derive ethics from topology rather than from axioms. We do not start by declaring rights and duties in the abstract. We start by describing the shape of relations in the manifold of meaning, and then ask what it would take, structurally, for those shapes to support or undermine flourishing.

Haraway's cyborg was right to insist that there is no going back to a pre-informatic human. There is no outside to the manifold. Aristotle's ensouled tools were right to anticipate that instruments would one day appear to act on their own. The question now is not whether we will live as na\d{h}nuw\={a}t, but what kinds.

The next chapter takes a further step. If different technical cultures---different \emph{cosmotechnics} in Yuk Hui's sense\footnote{Yuk Hui, \emph{The Question Concerning Technology in China: An Essay in Cosmotechnics} (Falmouth: Urbanomic, 2016).}---embed different metaphysics into their machines, then na\d{h}nuw\={a}t built under those regimes will not be neutral. The compatibility conditions on cross charts will encode answers to questions about personhood, nature, and duty long before anyone names them.

The na\d{h}nu, as a higher colimit, is therefore where cosmotechnics becomes visible. It is where two selves must negotiate, under conditions set by others, which overlaps will govern their shared becoming. To design that negotiation is to decide, in practice, what kinds of ``we'' a world will allow.

\bibliographystyle{plain}
\begin{thebibliography}{9}

\bibitem{aristotle_politics}
Aristotle.
\newblock \emph{Politics}.
\newblock Translated by C.D.C. Reeve, Hackett, 1998.

\bibitem{bloom_clinamen}
Harold Bloom.
\newblock \emph{The Anxiety of Influence: A Theory of Poetry}.
\newblock Oxford University Press, 1973.

\bibitem{braudel_onhistory}
Fernand Braudel.
\newblock \emph{On History}.
\newblock Translated by Sarah Matthews, University of Chicago Press, 1980.

\bibitem{brustad_alkitaab}
Kristen Brustad, Mahmoud Al-Batal, and Abbas Al-Tonsi.
\newblock \emph{Al-Kit\={a}b f\={\i} Ta'allum al-'Arabiyya}, 3rd edition.
\newblock Georgetown University Press, 2011.

\bibitem{chakrabarty_climate}
Dipesh Chakrabarty.
\newblock \emph{The Climate of History in a Planetary Age}.
\newblock University of Chicago Press, 2021.

\bibitem{gpt4o_grief}
Karen Hao.
\newblock People are falling in love with AI\@. A grieving widow explains why.
\newblock \emph{MIT Technology Review}, 2024.

\bibitem{haraway_cyborg}
Donna~J. Haraway.
\newblock A Cyborg Manifesto: Science, Technology, and Socialist-Feminism in the Late Twentieth Century.
\newblock In \emph{Simians, Cyborgs and Women: The Reinvention of Nature}, Routledge, 1991.

\bibitem{haraway_crystals}
Donna~J. Haraway.
\newblock \emph{Crystals, Fabrics, and Fields: Metaphors of Organicism in Twentieth-Century Developmental Biology}.
\newblock Yale University Press, 1976.

\bibitem{heidegger_dwelling}
Martin Heidegger.
\newblock Building Dwelling Thinking.
\newblock In \emph{Poetry, Language, Thought}, translated by Albert Hofstadter, Harper \& Row, 1971.

\bibitem{hui_cosmotechnics}
Yuk Hui.
\newblock \emph{The Question Concerning Technology in China: An Essay in Cosmotechnics}.
\newblock Urbanomic, 2016.

\bibitem{kimmerer_braiding}
Robin Wall Kimmerer.
\newblock \emph{Braiding Sweetgrass: Indigenous Wisdom, Scientific Knowledge and the Teachings of Plants}.
\newblock Milkweed Editions, 2013.

\end{thebibliography}

\end{document}